\documentclass[11pt]{article}

\usepackage[french, english]{babel}

\usepackage[T1]{fontenc}

\usepackage{amsfonts}
\usepackage{amsmath}
\usepackage{amssymb}

\usepackage{graphicx}
\usepackage{multirow}
\usepackage{multicol}
\usepackage{xcolor}
\usepackage[shortlabels]{enumitem}

\usepackage[margin=2cm]{geometry}

\usepackage{mathtools}
\usepackage{siunitx}
\usepackage{braket}

\usepackage{listings}

\usepackage{tikz}
\usetikzlibrary{shapes.geometric, arrows.meta, positioning}

\tikzstyle{block} = [rectangle, rounded corners, minimum width=4cm, minimum height=1cm,
text centered, draw=black, fill=blue!10]
\tikzstyle{process} = [rectangle, rounded corners, minimum width=5cm, minimum height=1cm,
text centered, draw=black, fill=green!10]
\tikzstyle{io} = [trapezium, trapezium left angle=70, trapezium right angle=110,
minimum width=4cm, minimum height=1cm, text centered, draw=black, fill=orange!15]
\tikzstyle{decision} = [diamond, aspect=2, draw=black, fill=yellow!15, text centered]
\tikzstyle{arrow} = [thick, ->, >=Stealth]


\renewcommand\lstlistingname{Code Python}

\lstset{frame=tb,
  language=Python,
  numbers=left,
  stepnumber=1,
  aboveskip=3mm,
  belowskip=3mm,
  showstringspaces=false,
  columns=flexible,
  basicstyle={\small\ttfamily},
%   numbers=none,
  numberstyle=\tiny\color{gray},
  keywordstyle=\color{blue},
  commentstyle=\color{paired2b},
  stringstyle=\color{paired3b},
  breaklines=true,
  breakatwhitespace=true,
  tabsize=1,
  frame=single,
%   showtabs=false,
}


\begin{document}


\hrule

\begin{center}
  {\Large GEI299 - Gestion de projet} \\ \vspace{5mm}
  {\LARGE\sffamily Devoir 2} \\
  Alexis Jean \\
  À remettre le 23 octobre
\end{center}

\hrule


\section{\underline{Introduction}}
L'objectif de ce document est de décrire le cahier de conception du logiciel à l'étude. Voyons un petit rappel de notre contrat : GE Healthcare nous demande de de créer un logiciel de simulation quantique pour pouvoir trouver un métal moins toxique que le gadolinium pour les radiographies. Le logiciel doit permettre à l'utilisateur de choisir entre deux modes : un mode scientifique pour les chercheurs et un mode affaire pour les décideurs. Le logiciel doit se connecter à la plateforme IBM Quantum pour effectuer les simulations et doit inclure un module d'analyse multi-critères pour évaluer les résultats en fonction de divers paramètres tels que le coût, la durabilité et la performance.
\paragraph{}
Rappelons aussi la portée du projet où nous devons livrer une version prototype ayant la possibilité d'être améliroré, donc que le prototype soit fait de façon modulaire. Égalment, il faut que le logiciel puisse être multiplateforme soit Windows et MacOS. Nous avons aussi des contraintes telles que le temps de calcul doi être inférieur à 2 minutes, le taux d'erreur par rapport à la référence classique doit être inférieur à 2\% et le contrat offre un budget maximal de 200k\$ pour le développement du logiciel.

\section{\underline{Diagrammes de la conception}}

\subsection{Diagramme conceptuel}
Voici ici le diagramme conceptuel de notre logiciel de simulation quantique :

\begin{figure}[htbp]
\centering
\begin{tikzpicture}[node distance=1.8cm]

\node (user) [io] {Utilisateur (Scientifique / Affaire)};
\node (ui) [block, below of=user] {Interface Utilisateur Graphique};
\node (sim) [process, below of=ui] {Moteur de Simulation Quantique (IBM Quantum)};
\node (analysis) [block, below of=sim] {Module d’Analyse Multi-critères};
\node (export) [process, below of=analysis] {Visualisation et Exportation des Résultats};

\draw [arrow] (user) -- node[right]{Entrée du choix de mode} (ui);
\draw [arrow] (ui) -- node[right]{Entrée des données pour la simulation} (sim);
\draw [arrow] (sim) -- node[right]{Résultats de la simulation} (analysis);
\draw [arrow] (analysis) -- node[right]{Classement et évaluation} (export);

\end{tikzpicture}
\caption{Diagramme conceptuel du logiciel de simulation quantique.}
\end{figure}

\newpage

\textbf{Légende :}
\begin{itemize}
    \item Trapezium orange : Utilisateur
    \item Bloc bleu : interface ou traitement interne du système
    \item Bloc vert : processus de calcul ou module logiciel
    \item Flèches : flux de données ou d’information
\end{itemize}

Maintenant, nous pouvons faire une brève description des différents composants du diagramme conceptuel :
\begin{itemize}
  \item \textbf{Utilisateur (Scientifique / Affaire)} : L'utilisateur peut choisir entre deux modes d'utilisation du logiciel, soit le mode scientifique pour les chercheurs ou le mode affaire pour les décideurs.
  \item \textbf{Interface Utilisateur Graphique} : Permet à l'utilisateur de sélectionner le mode d'utilisation et d'entrer les données nécessaires pour la simulation.
  \item \textbf{Moteur de Simulation Quantique (IBM Quantum)} : Se connecte à la plateforme IBM Quantum pour exécuter les simulations basées sur les données fournies par l'utilisateur.
  \item \textbf{Module d’Analyse Multi-critères} : Analyse les résultats de la simulation en fonction de divers critères tels que le coût, la durabilité et la performance et classe les différents candidats dans une base de données.
  \item \textbf{Visualisation et Exportation des Résultats} : Présente les résultats de l'analyse à l'utilisateur de manière claire et permet l'exportation des données pour une utilisation ultérieure.
\end{itemize}


\subsection{Schéma bloc}
Voici ici le schéma bloc de notre logiciel de simulation quantique :

\begin{figure}[h!]
\centering
\begin{tikzpicture}[node distance=1.8cm]

\node (ui) [block] {Interface Utilisateur};
\node (proc) [process, below of=ui] {Module de Traitement et Connexion IBM Quantum};
\node (db) [block, below of=proc] {Base de Données Interne};
\node (analysis) [process, below of=db] {Analyse Multi-Critères};
\node (export) [block, below of=analysis] {Module d’Exportation et Visualisation};

\draw [arrow] (ui) -- node[right]{Entrées utilisateur} (proc);
\draw [arrow] (proc) -- node[right]{Résultats bruts de simulation} (db);
\draw [arrow] (db) -- node[right]{Données structurées} (analysis);
\draw [arrow] (analysis) -- node[right]{Scores et classements} (export);

\end{tikzpicture}
\caption{Schéma bloc fonctionnel du système logiciel.}
\end{figure}

\newpage

\textbf{Légende :}
\begin{itemize}
    \item Bloc bleu : composant logiciel principal
    \item Bloc vert : module de calcul ou de traitement
    \item Flèches : circulation des données entre modules
\end{itemize}
En détail, chaque blocs de ce schéma bloc représente respectivement:

\begin{itemize}
  \item \textbf{Interface utilisateur} : L'utilisateur peux choisir le mode d'utilisation soit scientifique et affaires selon ses besoins. À partir de ce dernier, il peut choisir toutes les données à entrer pour la simulation, accéder à la base de données ainsi qu'au classement, etc.
  \item \textbf{Module de traitement de connexion IBM Quantum} : Ici, 'a partir des données qui ont été entrées par l'utilisateur, le logiciel va pouvoir traiter les informations, pouvoir préparer une simulation quantique et l'envoyer à la plateforme IBM Quantum pour l'exécution de cette dernière ainsi que la récupération des résultats. 
  \item \textbf{Base de données interne} : Le rôle de la base de données est d'enregister les résultats bruts qui été récupérés de la platerforme IBM Quantum et de les structurer pour qu'ils soient analysés par la suite dans le prochain module.
  \item \textbf{Analyse multi-critères} : L'analyse multi-critères va prendre les données structurées de la base de données et va évaluer les différents candidats en fonction des critères tels que le coût, la durabilité et la performance. Ensuite, elle va générer des scores et des classements pour chaque candidat.
  \item \textbf{Module d'exportation et visualisation} : Ce dernier module sert à exporter les résultats du classement fait par l'analyse multi-critères sous les formats désirés par l'utilisateur (ex : PDF, Jupyter Notebook, etc.) et aussi de visualiser les résultats de manière claire et compréhensible pour l'utilisateur.
\end{itemize}

\newpage

\section{\underline{Décomposition du schéma bloc}}
On peut décomposer les blocs du schéma bloc en faisant un tableau de description des modules :
\begin{table}[h!]
\centering
\caption{Modules fonctionnels du système de simulation quantique}
\begin{tabular}{|p{3cm}|p{4cm}|p{3cm}|p{3cm}|p{2cm}|}
\hline
\textbf{Module} & \textbf{Fonction principale} & \textbf{Entrées} & \textbf{Sorties} & \textbf{Liens} \\
\hline
Interface utilisateur & Gérer les interactions avec l’usager (sélection de métal, ligand, mode) & Commandes et paramètres de l’utilisateur (clavier)& Données de simulation formatées & Moteur de simulation \\
\hline
Moteur de simulation quantique & Préparer une simulation quantique et envoyer les calculs sur la plateforme IBM Quantum et récupérer les résultats & Données du métal et du ligand & Résultats quantiques (énergie fondamentale, propriétés électroniques) & Base de données interne \\
\hline
Base de données interne & Stocker les métaux, ligands et résultats des simulations & Données brutes issues de la simulation de IBM Quantum & Données consultables et réutilisables & Analyse multi-critères \\
\hline
Analyse multi-critères & Évaluer et comparer les métaux selon plusieurs critères (relaxivité, toxicité, coût) & Données de la base de données & Classement et indicateurs de performance & Module d’exportation \\
\hline
Module d’exportation et visualisation & Générer des rapports, graphiques et exports (PDF, Jupyter notebook, etc.) & Données d’analyse & Fichiers et affichages pour l’utilisateur & Interface utilisateur \\
\hline
\end{tabular}
\end{table}

\newpage

\section{\underline{Spécifications fonctionnelles et non-fonctionnelles}}

\subsection{Spécifications fonctionnelles}
Dans cette section, nous allons détailler les spécifications fonctionnelles du logiciel de simulation quantique. Ces spécifications décrivent les fonctionnalités que le logiciel doit offrir pour répondre aux besoins des utilisateurs. Le tout sera en forme de tableau : 

\begin{table}[h!]
\centering
\caption{Spécifications fonctionnelles du système}
\begin{tabular}{|p{1cm}|p{12cm}|}
\hline
\textbf{ID} & \textbf{Fonction attendue} \\
\hline
F1 & Le système doit permettre d’ajouter, modifier et supprimer des métaux candidats. \\
\hline
F2 & L’utilisateur doit pouvoir sélectionner un ligand associé et configurer les paramètres de simulation. \\
\hline
F3 & Le logiciel doit établir une connexion sécurisée à la plateforme IBM Quantum pour lancer les calculs. \\
\hline
F4 & Les résultats des simulations doivent être enregistrés automatiquement dans la base de données interne. \\
\hline
F5 & Un module doit analyser les résultats selon les critères : relaxivité, toxicité, biodégradabilité, coût. \\
\hline
F6 & Le système doit permettre d’exporter les résultats en formats PDF, Jupyter notebook, etc. \\
\hline
F7 & L’utilisateur doit pouvoir basculer entre un mode “scientifique” et un mode “affaires”. \\
\hline
\end{tabular}
\end{table}

\subsection{Spécifications non-fonctionnelles}
Maintenant, voici les spécifications non-fonctionnelles du logiciel de simulation quantique. Ces spécifications définissent les critères de performance, de sécurité et d’autres aspects qualitatifs du système. Le tout sera aussi en forme de tableau :
\begin{table}[h!]
\centering
\caption{Spécifications non-fonctionnelles du système}
\begin{tabular}{|p{1cm}|p{5cm}|p{8cm}|}
\hline
\textbf{ID} & \textbf{Catégorie} & \textbf{Description} \\
\hline
NF1 & Performance & Le temps de calcul d’une simulation standard doit être inférieur à 2 minutes. \\
\hline
NF2 & Précision & L’erreur moyenne par rapport à une référence classique doit être inférieure à 2\%. \\
\hline
NF3 & Compatibilité & Le logiciel doit fonctionner sous Windows et MacOS. \\
\hline
NF4 & Sécurité & Les communications avec IBM Quantum doivent être chiffrées. \\
\hline
NF5 & Ergonomie & L’interface doit être intuitive, accessible et bilingue (FR/EN). \\
\hline
NF6 & Évolutivité & Le code doit être modulaire pour permettre l’ajout de nouveaux algorithmes. \\
\hline
NF7 & Traçabilité & Les simulations doivent être automatiquement journalisées. \\
\hline
\end{tabular}
\end{table}


\section{\underline{Objectfs SMART}}
On applique la méthode SMART pour définir des objectifs clairs et atteignables pour le projet de développement du logiciel de simulation quantique. Voici les objectifs SMART pour ce projet :

\begin{itemize}
  \item \textbf{Livraison du prototype} : Le logiciel doit être fonctionnel et livré au client dans un délai de 2 mois.
  \item \textbf{Performance de simulation} : Le temps de calcul pour une simulation standard doit être inférieur à 2 minutes, avec une précision de résultats inférieure à 2\% d'erreur par rapport aux références classiques.
  \item \textbf{Maintenance et évolutivité} : Le code doit être modulaire et bien documenté pour permettre des mises à jour et l'ajout de nouvelles fonctionnalités dans les 6 mois suivant la livraison initiale.
\end{itemize}

\section{\underline{Livrables}}
Dans cette section, les livrables du projet seront brièvement énumérés selon les catégorie établies :

\begin{itemize}
  \item \textbf{Livrables techniques} :
    \begin{itemize}
      \item Un prototype fonctionnel du logiciel de simulation quantique.
      \item Le code source de ce prototype pour les améliorations futures.
      \item La base de données interne avec les résultats des simulations.
    \end{itemize}
  \item \textbf{Livrables documentaires} :
    \begin{itemize}
      \item Le rapport technique du résultat de chaque simulation et comparaison.
      \item Le manuel utilisateur pour les deux modes (scientifique et affaires).
    \end{itemize}
  \item \textbf{Livrables organisationnels} :
    \begin{itemize}
      \item Une présentation des résulats pour les affaires.
      \item Le cahier Jupyter Notebook avec les analyses et visualisations des résultats.
    \end{itemize}
\end{itemize}


\section{\underline{Validation et tests}}
Cette section décrit les différentes méthodes de validation ainsi que les tests qui serviront à vérifier la conformité du logiciel aux spécifications établies. Le tout est décrit dans un tableau à la page suivante :

\begin{table}[h!]
\centering
\caption{Plan de validation et de tests du système}
\begin{tabular}{|p{1cm}|p{4cm}|p{5cm}|p{3cm}|p{3cm}|}
\hline
\textbf{ID} & \textbf{Élément testé} & \textbf{Critère de validation / Objectif du test} & \textbf{Méthode de test} & \textbf{Résultat attendu} \\
\hline
T1 & Interface utilisateur & Vérifier que toutes les commandes (ajout, simulation, export) sont accessibles et fonctionnelles. & Test manuel d’interaction / scénario utilisateur. & Aucune erreur, navigation fluide, interface intuitive. \\
\hline
T2 & Connexion IBM Quantum & Vérifier l’établissement d’une connexion sécurisée et la stabilité de la communication. & Test d’intégration avec API IBM / journal réseau. & Connexion stable, chiffrement actif, aucune perte de paquets. \\
\hline
T3 & Module de simulation quantique & Confirmer que les résultats calculés (énergie fondamentale, propriétés) sont cohérents avec les valeurs classiques de référence. & Comparaison avec un benchmark de simulation classique. & Écart maximal < 2\% pour les propriétés clés. \\
\hline
T4 & Base de données interne & Vérifier la sauvegarde, la mise à jour et la consultation correcte des résultats de simulation. & Tests unitaires d’écriture/lecture et scénarios utilisateurs multiples. & Toutes les données sont enregistrées et récupérées sans corruption. \\
\hline
T5 & Module d’analyse multi-critères & Confirmer que les scores et classements sont calculés correctement à partir des critères définis. & Vérification algorithmique et comparaison avec jeu de données de test connu. & Classements cohérents avec les poids de critères fixés. \\
\hline
T6 & Module d’exportation & Vérifier la génération correcte des fichiers (PDF, Jupyter Notebook, etc.) et leur intégrité. & Test de sortie automatique et ouverture des fichiers générés. & Fichiers lisibles, formats conformes et contenus exacts. \\
\hline
T7 & Performance du système & Mesurer le temps d’exécution global d’une simulation et la consommation de ressources. & Profilage de performance sur plusieurs scénarios. & Temps total < 2 min, utilisation CPU/mémoire acceptable. \\
\hline
T8 & Sécurité et journalisation & Vérifier que toutes les simulations sont tracées et qu’aucune donnée sensible n’est exposée. & Audit des journaux et tentative d’accès non autorisé simulée. & Journalisation complète, accès non autorisés bloqués. \\
\hline
T9 & Compatibilité multiplateforme & Tester le fonctionnement sous Windows et macOS. & Tests d’installation et d’exécution sur les deux OS. & Logiciel pleinement fonctionnel sur les deux plateformes. \\
\hline
T10 & Ergonomie / accessibilité & Vérifier la lisibilité, le contraste, et la navigation au clavier/souris. & Tests d’utilisabilité avec échantillon d’utilisateurs. & 90\% des utilisateurs évaluent l’ergonomie comme satisfaisante ou excellente. \\
\hline
\end{tabular}
\end{table}



\end{document}







